\documentclass[11pt]{article}
\usepackage[margin=1in]{geometry}
\usepackage{amsmath, amssymb, amsthm}
\usepackage{booktabs}

\title{Problem 463: A Weird Recurrence Relation}
\author{}
\date{}

\newtheorem{theorem}{Theorem}
\newtheorem{lemma}[theorem]{Lemma}

\begin{document}
\maketitle

\section{Problem Statement}

Define $f$ on positive integers:
\begin{align*}
f(1) &= 1,\quad f(3) = 3, \\
f(2n) &= f(n), \\
f(4n+1) &= 2f(2n+1) - f(n), \\
f(4n+3) &= 3f(2n+1) - 2f(n).
\end{align*}
Let $S(n) = \sum_{i=1}^{n} f(i)$. Given $S(8) = 22$, $S(100) = 3604$, find $S(3^{37}) \bmod 10^9$.

\section{Mathematical Foundation}

\begin{lemma}[Dependence on odd part]
\label{lem:odd}
For all positive integers $n$, $f(n) = f(\mathrm{odd}(n))$, where $\mathrm{odd}(n)$ is the largest odd divisor of~$n$.
\end{lemma}

\begin{proof}
Write $n = 2^s m$ with $m$ odd. We induct on~$s$. The base $s = 0$ is trivial. For $s \ge 1$, $f(n) = f(2 \cdot 2^{s-1}m) = f(2^{s-1}m)$ by the rule $f(2n) = f(n)$. By induction, $f(2^{s-1}m) = f(m)$.
\end{proof}

\begin{theorem}[Even--odd splitting]
\label{thm:split}
Define $T(k) = \sum_{j=0}^{k} f(2j+1)$. Then:
\begin{align}
S(2N) &= S(N) + T(N-1), \label{eq:Seven} \\
S(2N+1) &= S(N) + T(N). \label{eq:Sodd}
\end{align}
\end{theorem}

\begin{proof}
Split the sum $S(2N) = \sum_{i=1}^{2N} f(i)$ by parity of~$i$:
\[
S(2N) = \sum_{i=1}^{N} f(2i) + \sum_{i=1}^{N} f(2i-1) = \sum_{i=1}^{N} f(i) + \sum_{j=0}^{N-1} f(2j+1) = S(N) + T(N-1),
\]
using Lemma~\ref{lem:odd} for the even-indexed terms. Adding $f(2N+1)$ yields~\eqref{eq:Sodd}.
\end{proof}

\begin{theorem}[Recurrence for $T$]
\label{thm:T}
For $k \ge 2$:
\begin{align}
T(2k) &= 2T(k) + 3T(k-1) - S(k) - 2S(k-1) - 1, \label{eq:Teven} \\
T(2k+1) &= T(2k) + 3f(2k+1) - 2f(k). \label{eq:Todd}
\end{align}
Base cases: $T(0) = 1$, $T(1) = 4$.
\end{theorem}

\begin{proof}
Partition the odd indices $\{1, 3, \ldots, 4k+1\}$ by residue modulo~4. Substituting the recurrences for $f(4j+1)$ and $f(4j+3)$ into the respective partial sums and regrouping in terms of $T$ and $S$ yields~\eqref{eq:Teven}. Equation~\eqref{eq:Todd} follows directly from $T(2k+1) = T(2k) + f(4k+3) = T(2k) + 3f(2k+1) - 2f(k)$.
\end{proof}

\begin{theorem}[Complexity]
\label{thm:complexity}
The mutual recursion $(S, T, f)$ with memoization computes $S(n) \bmod 10^9$ in $O(\log^2 n)$ time and $O(\log^2 n)$ space.
\end{theorem}

\begin{proof}
Each recursive call halves its argument. Starting from~$n$, the distinct arguments form a set of size $O(\log n)$ per recursion level. Cross-calls between $S$, $T$, and $f$ introduce $O(\log n)$ new arguments per level, yielding $O(\log^2 n)$ total subproblems. Each requires $O(1)$ modular arithmetic operations.
\end{proof}

\section{Algorithm}

The algorithm implements the mutual recursion $(S, T, f)$ with hash-map memoization:
\begin{enumerate}
    \item Compute $f(n)$ by reducing even arguments ($f(2n) = f(n)$) and applying the mod-4 recurrence for odd arguments.
    \item Compute $T(k)$ via Theorem~\ref{thm:T}, calling $S$ and $f$ as needed.
    \item Compute $S(n)$ via Theorem~\ref{thm:split}, calling $T$ as needed.
    \item All arithmetic is performed modulo $10^9$.
\end{enumerate}

\section{Complexity Analysis}

\begin{itemize}
    \item \textbf{Time:} $O(\log^2 n)$. For $n = 3^{37} \approx 4.5 \times 10^{17}$, $\log_2 n \approx 59$, giving ${\sim}3500$ subproblems.
    \item \textbf{Space:} $O(\log^2 n)$ for the memoization tables.
\end{itemize}

\section{Verification}

\begin{center}
\begin{tabular}{cr}
\toprule
$n$ & $S(n)$ \\
\midrule
1 & 1 \\
8 & 22 \\
100 & 3604 \\
\bottomrule
\end{tabular}
\end{center}

\section{Answer}

\[
\boxed{808981553}
\]

\end{document}
